\scanpage{115}
\mspage{106}
% archive sheet number: 115 (lower left)
% L1-4
Reste le cas $\deleted{\unsure{Z'>}}P_a\inserted{<Z'}$. Mais comme $Z<X$, on aura $Z\leq P_a$ ou \deleted{\unsure{$Z\leq P_a$|$Z<P_a$}}, donc on aura $Z<Z'$, d'où encore $Z\mathrel{\mid\circ\mid}Z'$, qed

% L5-8
\uline{NB} Je n'ai pas eu à utiliser Hyp., \deleted{ni la \unsure{constructibilité} de $\Phi$} mais \deleted{la} j'ai utilisé \deleted{\unsure{la const.|le contre}} que $\Phi$ $\Psi$ sont des \unsure{préfigures}.

% L9-14
\uline{Dém de (b)} Résulte par récurrence \inserted{(sur $\operatorname{card}\delta\Phi$)} de $1^{\circ}$, et du \deleted{\unsure{Deux}}

\uline{Lemme} : \unsure{Soient|Soit} $a<b<c$, et $\Phi=\{S_{ab},P_b,S_{bc}\}$. Alors
\[
\operatorname{supp}^{\circ}(\Phi)=\operatorname{supp}^{\circ}(S_{ac})
\]
(Lemme de composition des segments).

% L15-20
\uline{Démonstration} Il suffit de voir que $\Phi$ et $S_{ac}$ ont \unsure{même|mêmes} $\operatorname{cosupp}^{\circ}$. \deleted{Soit} Cela résulte du résultat légèrement plus \deleted{\unsure{donc} $X$} précis :

% L21-23
Si $X\in\unsure{M|\mathcal{M}}$, on a
\[
\begin{aligned}
X<_{\mathrm{pos}}\Phi\quad\inserted{\text{(i.e. }X<S_{ab}\text{)}}&\Longleftrightarrow X<S_{ac}\\
X>_{\mathrm{pos}}\Phi\quad\inserted{\text{(i.e. }X>S_{bc}\text{)}}&\Longleftrightarrow X>S_{ac}
\end{aligned}
\]
\margin{trivial (\unsure{sans Hyp.})}
